\chapter{Revisão Teórica}
\label{cap:revisao}

\section{Realidade Virtual na Educação}

O estudo de \citeonline{merchant2014} realizou uma meta-análise para avaliar a eficácia da instrução baseada em realidade virtual (RV) no ensino fundamental, médio e superior. Foram analisados 67 estudos categorizados em três formatos de RV: jogos, simulações e mundos virtuais. Os resultados indicaram que todos os formatos de RV contribuíram para a melhoria dos resultados de aprendizagem, mas os jogos mostraram maior impacto em comparação com simulações e mundos virtuais \cite[p.~30]{merchant2014}. Essa vantagem foi atribuída à estrutura interativa dos jogos, que promovem maior engajamento e motivação dos alunos. Além disso, os jogos permitem estratégias de resolução de problemas e testes de hipóteses, elementos fundamentais para o aprendizado ativo.

A pesquisa também revelou que a retenção do conhecimento a longo prazo foi mais eficaz nos ambientes gamificados, enquanto o impacto das simulações e dos mundos virtuais foi mais dependente das condições de teste. Para as simulações, os melhores resultados foram observados quando utilizadas como complemento ao ensino tradicional, permitindo aos alunos praticarem conceitos previamente introduzidos. Já os mundos virtuais, apesar de proporcionarem um ambiente aberto para interação e construção de conhecimento, apresentaram menor efetividade quando medidos ao longo do tempo. O estudo destaca a importância do \textit{design} instrucional adequado para maximizar o potencial da RV na educação, enfatizando a necessidade de estratégias pedagógicas que combinem os benefícios das diferentes abordagens de aprendizado imersivo \cite[p.~37]{merchant2014}.

\section{Aplicabilidade da Realidade Virtual no Treinamento em Segurança na Construção Civil}

A utilização da Realidade Virtual (RV) no treinamento em segurança na construção civil tem sido amplamente discutida na literatura devido ao seu potencial de oferecer um ambiente imersivo e interativo para a capacitação dos trabalhadores. \citeonline[p.~9-10]{akindele2024} destacam que a RV permite uma experiência de aprendizagem envolvente e realista, contribuindo para que os trabalhadores compreendam de forma mais adequada os riscos e adotem estratégias apropriadas para lidar com diferentes situações de perigo nos canteiros de obras. Essa tecnologia permite criar ambientes simulados para identificação de riscos, treinamento operacional e educação em segurança \cite{zhou2019,wang2020,akindele2024}.

No que diz respeito à identificação de perigos, a RV permite que os trabalhadores interajam com ambientes virtuais projetados para replicar situações reais de risco. Segundo \citeonline[p.~10]{akindele2024}, estudos demonstram que a utilização da RV melhora a capacidade de identificação de riscos em relação a métodos tradicionais, como fotografias e documentação técnica, tornando a detecção mais precisa e eficaz. A tecnologia oferece um ambiente seguro para simulação de cenários de perigo, como trabalho em altura e manuseio de equipamentos, sem riscos reais \cite{li2018,sacks2019,akindele2024}. Estudos comparativos indicam que a simulação imersiva aumenta a retenção do conhecimento e a velocidade de resposta a situações de risco quando comparada a métodos convencionais, como manuais e apresentações \cite{bhoir2015,akindele2024}.

Outro aspecto abordado é o uso da RV para o treinamento de procedimentos. A tecnologia permite que os trabalhadores pratiquem atividades críticas antes de executá-las no ambiente real, reduzindo a probabilidade de erros e aumentando a eficiência \cite[p.~10]{akindele2024}. Essa abordagem é particularmente vantajosa para a preparação de respostas a situações de emergência, nas quais a rapidez e a precisão das ações podem ser determinantes para a segurança \cite{fang2016,guo2017,akindele2024}.

A RV também se mostra uma ferramenta eficaz para a educação em segurança, sendo mais envolvente que métodos tradicionais, como palestras \cite{sacks2013,guo2017}. \citeonline[p.~10]{akindele2024} ressaltam que, ao simular tarefas e situações de risco, a RV permite que os trabalhadores adquiram conhecimentos de forma intuitiva e de fácil retenção. Além disso, a interatividade e a imersão proporcionadas por essa tecnologia favorecem a retenção de informações e o desenvolvimento de boas práticas de segurança, mitigando desafios da aprendizagem passiva \cite{akindele2024}.

\subsection{Integração da RV com a Segurança do Trabalho}

\citeonline{oliveira2024} complementam essa visão, observando que, embora a RV esteja em expansão na construção civil, a \textit{integração explícita} com a segurança do trabalho ainda é limitada. Em sua revisão sistemática, apenas 25\% dos estudos abordavam diretamente a segurança. Isso sugere que, embora a RV seja usada para treinamento de forma geral, o \textit{foco específico} na segurança ainda precisa ser ampliado.

Os autores também destacam a importância da integração da RV com metodologias como o \gls{BIM}. Essa integração permite um melhor planejamento e visualização de projetos, o que contribui para a identificação proativa de riscos e a melhoria das condições de segurança nos canteiros de obras \cite[p.~12]{oliveira2024}.

\subsection{Desafios e Lacunas na Pesquisa}

Apesar das vantagens, diversos desafios e lacunas na pesquisa sobre a aplicação da RV no treinamento em segurança ainda precisam ser superados. \citeonline[p.~12-13]{akindele2024} apontam a escassez de estudos comparativos que avaliem o impacto da aprendizagem mediada por RV em relação aos métodos tradicionais. Além disso, destacam a necessidade de investigar como variáveis individuais, como experiência prévia e conhecimento sobre segurança, afetam a eficácia do treinamento virtual. Segundo os autores, é preciso investigar se ``a experiência prévia e o nível de compreensão de segurança impactam os resultados de aprendizado dentro de um ambiente virtual'' \cite[p.~13]{akindele2024}. Estudar como a carga cognitiva influencia o desempenho do treinamento em RV pode fornecer novas direções para pesquisas futuras \cite{zhao2015,lucas2018}.

Outro desafio é a personalização dos conteúdos de RV para diferentes perfis de usuários (trabalhadores, engenheiros, gestores). \citeonline{akindele2024} enfatizam que a RV deve ser adaptada às necessidades de cada grupo. Além disso, sugerem que a RV pode ter um papel importante na redução de riscos relacionados a fatores cognitivos, como distrações, mas ressaltam que essa abordagem ainda carece de investigação \cite[p.~13]{akindele2024}.

\subsection{Recomendações para Estudos Futuros}

\citeonline[p.~14]{akindele2024} recomendam estudos mais robustos, com grupos de controle, para comparar métodos de treinamento em RV. Além disso, sugerem investigar como a experiência prévia dos trabalhadores influencia a compreensão dos conceitos de segurança no ambiente virtual. Os autores afirmam que ``um método prático de teste com diversos grupos de aprendizes e uso de grupos de controle pode fornecer bons \textit{insights}'' \cite[p.~14]{akindele2024}. Estudos futuros também poderiam investigar a combinação da RV com outras tecnologias, como realidade aumentada e sensores biométricos. \citeonline{oliveira2024} enfatizam a necessidade de maior integração entre os sistemas imersivos e as \textit{normativas de segurança}, além do desenvolvimento de estratégias para viabilizar a implementação da RV de forma acessível e eficiente.

Em suma, a literatura evidencia o potencial da RV para transformar o treinamento em segurança na construção civil, mas é fundamental que estudos futuros abordem as limitações identificadas, garantindo que a tecnologia seja utilizada de maneira eficaz e adaptada às necessidades dos profissionais da construção.

\section{Métodos de Avaliação}

Com base dos artigos selecionados, é possível fazer observações de como a efetividade dessas experiências em RV foram avaliadas por outros pesquisadores:

\subsection{Abotaleb et al. (2023)}

\citeonline{abotaleb2023} avaliaram a eficácia do treinamento em RV em comparação com um método tradicional baseado no curso ``IOSH Managing Safely\textregistered''. Para isso, estudantes de engenharia foram divididos aleatoriamente em dois grupos: um que recebeu treinamento tradicional e outro que utilizou a simulação interativa de um canteiro de obras em RV. A avaliação foi realizada por meio de um teste padronizado após o treinamento, que media a capacidade dos participantes de identificar perigos e aplicar conhecimentos de segurança. Os resultados foram analisados estatisticamente usando testes de normalidade (Shapiro-Wilk), comparação de médias (teste t ou Mann-Whitney U, dependendo da distribuição dos dados) e o teste de Levene para verificar homogeneidade de variâncias. Os achados demonstraram um desempenho significativamente superior do grupo treinado com RV, e a diferença foi confirmada estatisticamente ($p = 0{,}01$), validando a efetividade do método imersivo de aprendizagem.

\subsection{Nykänen et al. (2020)}

\citeonline{nykanen2020} adotaram um ensaio controlado randomizado (\gls{RCT}) para comparar a eficácia do treinamento em RV com o treinamento tradicional baseado em palestras. Os participantes foram divididos em grupos, e diversas medidas foram coletadas antes, imediatamente após e um mês depois do treinamento. As avaliações incluíram a autoeficácia em segurança, o \textit{locus} de controle de segurança, o conhecimento sobre identificação de perigos e a motivação para comportamentos seguros. A experiência no ambiente RV foi analisada por meio de questionários que mediram presença, envolvimento, realismo e controle, além da incidência de sintomas de simulador. Dados observacionais, como padrões de movimento dentro do ambiente virtual, foram registrados e entrevistas qualitativas foram conduzidas para capturar percepções mais detalhadas. A análise estatística utilizou modelos lineares mistos generalizados (\gls{GLMM}) para examinar as diferenças entre os grupos ao longo do tempo, garantindo uma avaliação robusta da eficácia do treinamento.

\subsection{Xu e Zheng (2021)}

\citeonline{xu2021} avaliaram a efetividade do treinamento em realidade virtual (RV) para segurança em canteiros de obras através de um experimento comparativo, onde os participantes passaram tanto pelo treinamento tradicional (vídeo instrucional) quanto pelo treinamento em RV. A avaliação foi conduzida por meio de questionários antes e depois das sessões, medindo percepções sobre a plataforma RV, realismo da simulação e desconforto do simulador. Além disso, a análise do desempenho no ambiente virtual inclui métricas quantitativas, como a taxa de erro dos participantes e o Índice de Identificação de Perigo (\gls{HII}), além da observação direta do comportamento dos trabalhadores no ambiente RV. Para complementar os dados, questionários adicionais foram aplicados após a experiência em RV, incluindo a Escala de Aceitação de Van der Laan, e análises estatísticas foram conduzidas para verificar correlações entre o nível de dificuldade dos cenários e o desempenho dos participantes.

\subsection{Joshi et al. (2021)}

Na pesquisa de \citeonline{joshi2021}, a avaliação da eficácia e efetividade do treinamento em RV para segurança em canteiros de obras foi realizada por meio de um experimento comparativo entre dois grupos de participantes: um grupo que recebeu treinamento tradicional por vídeo e outro que utilizou um módulo RV interativo. A eficácia do sistema RV foi mensurada utilizando três métricas principais: \gls{SSQ} \cite{kennedy1993}, para avaliar o nível de desconforto causado pela simulação; SUS \cite{brooke1996}, para medir a usabilidade do sistema; e \gls{PQ}, para avaliar a experiência do usuário e a sensação de imersão no ambiente virtual. Já a efetividade do treinamento foi analisada com base em duas variáveis: motivação dos participantes, medida por meio de questionários comparando os dois métodos de treinamento, e ganho de conhecimento, avaliado por meio de testes aplicados antes e depois do treinamento. Análises estatísticas, como testes $t$ e qui-quadrado, foram utilizadas para verificar a significância das diferenças entre os grupos. Os resultados indicaram que o treinamento em RV gerou um maior nível de motivação nos participantes em comparação ao método tradicional, além de demonstrar potencial para melhorar a retenção de conhecimento sobre segurança no trabalho.

\section{Gaps Identificados}

Através dos artigos lidos, os problemas mais identificados foram:

\textbf{Integração insuficiente com segurança laboral:}
\begin{itemize}
    \item Apenas uma minoria dos estudos aborda explicitamente a correlação entre RV/RA e práticas de segurança, com poucos exemplos práticos de implementação em canteiros de obras reais.
\end{itemize}

\textbf{Limitações metodológicas:}
\begin{itemize}
    \item Amostras pequenas, falta de diversidade (ex.: trabalhadores experientes) e ausência de métricas quantitativas robustas comprometem a generalização dos resultados.
\end{itemize}

\textbf{Desafios na Implementação de Tecnologias de Segurança na Construção Civil:}
\begin{itemize}
    \item \textbf{Falta de Validação Prática:} Ausência de dados empíricos de longo prazo sobre a eficácia das tecnologias na redução de acidentes e custos, além de dúvidas sobre sua viabilidade econômica e adaptação a cenários complexos.
    \item \textbf{Barreiras de Comunicação e Implementação:} A falta de alinhamento entre academia e indústria resulta em tecnologias pouco adaptadas às necessidades práticas, com desafios como interoperabilidade, custos elevados e resistência à mudança em um setor pressionado por produtividade e orçamento.
    \item \textbf{Falta de Colaboração Multissetorial:} A escassez de parcerias entre pesquisadores, empresas e trabalhadores dificulta a implementação em escala real, enquanto o desconhecimento ou a falta de recursos perpetuam o uso de métodos tradicionais.
\end{itemize}

\textbf{Falta de Exemplos no Contexto Nacional:}
\begin{itemize}
    \item A NR-18 \cite{brasil2025nr18} estabelece requisitos específicos para segurança, como uso de EPIs, sinalização e prevenção de quedas. Contudo, protótipos existentes de RV (majoritariamente desenvolvidos no exterior) não incorporam esses padrões brasileiros, limitando sua aplicabilidade local. Não foi encontrado nenhum estudo que teste a eficácia de treinamentos em RV \textbf{sob as NRs}, especialmente em cenários reais de canteiros de obras brasileiros. Não há métricas claras sobre redução de acidentes ou conformidade normativa após o uso dessas ferramentas.
    \item Protótipos internacionais não consideram particularidades da construção civil brasileira, como práticas informais, diversidade de obras (desde grandes empreendimentos urbanos a construções rurais) e desafios socioeconômicos (ex.: capacitação de mão de obra menos especializada).
\end{itemize}
